% 华为杯全国研究生数学建模竞赛模板
% 基于 gmcmthesis 文档类
\documentclass[bwprint]{gmcmthesis}

% === 额外宏包（cls 未内置的）===
\usepackage[framemethod=TikZ]{mdframed}
\usepackage[section]{placeins}
\usepackage{needspace}
\usepackage{subfig}
\usepackage{float}
\usepackage{algorithm2e}

% === TikZ ===
\usepackage{tikz}
\usetikzlibrary{arrows.meta,positioning,shapes.geometric,fit,calc}
\usepackage{longtable}

% === 引用上标格式 ===
\setcitestyle{super,square,comma}

% ⛔ 摘要后空白页修复：已迁移到 gmcmthesis.cls（abstract 结尾用 \clearpage
%    替代 \newpage\null，避免内容接近页底时挤出含 \null 的空白页）。
%    此处不需要再 \renewenvironment{abstract}。

% === 竞赛信息 ===
\title{[论文标题]}
\baominghao{[参赛队号]}
\schoolname{[学校名称]}
\membera{[成员A]}
\memberb{[成员B]}
\memberc{[成员C]}

\begin{document}

\maketitle

% === 摘要 ===
\begin{abstract}

[中文摘要内容：华为杯默认走「丰满模式」——总字数 1500--2200 字、允许并鼓励跨两页；每个子问题独立一段（每段 280--400 字），段内按"过程式叙述 + 候选方法对比 + 中间过程数值 + 可视化手段"展开，并用 \textbf{} 各突出 1--3 处关键方法/结果。通用 600--800 字单页规则在华为杯不适用]

\keywords{[关键词1]\quad [关键词2]\quad [关键词3]\quad ... 华为杯默认 8--12 个（算法名+数学工具+模型族全列）；关键词行过长时用 \\ 在词间主动断行，禁止让某个关键词被拆到两行}
\end{abstract}

\pagestyle{plain}

% === 摘要与正文之间的换页（竞赛模板不出目录）===
% ⛔ 用 \clearpage 而非 \newpage：cls 修复后 abstract 已用 \clearpage 换页，
%    这里再用 \clearpage 检测到已在新页开头会自动跳过，不会产生空白页。
\clearpage
% ⛔ 不要在这里加回 \tableofcontents：竞赛论文统一不出目录页。
%   旁边的 \newpage/\clearpage 是摘要与正文之间的换页，删了正文会挤上来。
\clearpage

% === 正文 ===
\input{sections/1_restatement}
\input{sections/2_analysis}
\input{sections/3_assumptions}
\input{sections/4_symbols}
\input{sections/5_problem1}
\input{sections/6_problem2}
\input{sections/7_problem3}
\input{sections/8_sensitivity}
\input{sections/9_evaluation}

% === 参考文献 ===
% === AI 工具使用声明（自动插入）===
% 由 build_ai_disclosure.py 根据用户确认记录生成，请勿由模型改写。
\section*{AI 工具使用声明}
本参赛队在竞赛过程中使用了AI工具，主要用于语言润色、排版检查、参考文献格式整理，详细使用情况见支撑材料。

\begin{thebibliography}{99}
% \bibitem[1]{ref1} 作者. 标题[J]. 期刊, 年份.
\end{thebibliography}

% === 附录 ===
\begin{appendices}
\input{sections/A_code}
\end{appendices}

\end{document}
